\subsection{A smaller weighted moment suffices for generic returns}
\label{sec:flag-generic-reduced-moment}

Retain the original equal-word generic data of
Section~\ref{sec:flag-generic-pinned-input}, including its
actual absolutely convergent return expansion and finite
class and normal dictionary. We give a direct replacement
for the two Cauchy--Schwarz brackets in the generic
arithmetic assembly; in particular we do not enlarge
the mass of the local witness to (GP4).

\begin{lemma}[Generic-only local second moment]
\label{lem:flag-generic-only-local}
In the notation of (GL1)--(GL7), for any occupied integral
reference pairing class \(C_*\) one has
\[
 \sum_{\delta\bmod H_J}\mathrm{wt}(\delta)|A_\delta^0|^2
 \le 2P^{2s^2+2s}\,ngen(C_*).
 \tag{RM1}
\]
The factor two can be replaced by one when \(P=1\).
\end{lemma}
\begin{proof}
By (GL4)--(GL5),
\(m_\delta^0\le P^{s(s-1)+4s}\).
Use \(|A_\delta^0|^2\le(m_\delta^0)^2\), but sum the
generic mass itself rather than the mass over all
saturations. The exact generic equality in (GL6) gives
\[
 \sum_{\delta\bmod H_J}\mathrm{wt}(\delta)m_\delta^0
   =\sum_C w(C){\cal M}(M(C))
   \le \tfrac54 c_x P^{s(s-1)} .
\]
All corrected bases have the same determinant scaling,
and \(ngen(C_*)\ge3c_x/4\) as in (GL7).
The exponent after multiplication is
\(2s(s-1)+4s=2s^2+2s\).
At \(P=1\) the unique-class calculation of (GL7)
gives equality. No contribution of a proper saturation
occurs in this calculation.
\end{proof}

\begin{lemma}[A field-uniform endpoint sum]
\label{lem:flag-endpoint-packing}
For a CM field \(E\) of degree \(2f\), put
\[
 W(z)=\prod_{\tau}(1+|\tau z|^2)^2 ,
\]
using one complex embedding above each real place.
Then
\[
 \sum_{z\in\mathcal O_E} W(z)^{-1}
        \le (100/9)^f < 12^f .
\]
On the original space of block points of size \(s\),
restricted to nonzero block-end minors in \(\mathcal O_E\),
one consequently has
\[
 \sum_w \mathrm{blockMinorWeight}(w)^{-1}
       \le 14^{fs}.                                      \tag{RM2}
\]
\end{lemma}
\begin{proof}
The Minkowski image of two distinct algebraic integers
has difference of complex supremum norm at least one,
by the nonzero integer norm. The open products of disks
of radius \(1/2\) about their images are disjoint.
If \(|u-z|\le1/2\), then
\[
 1+|u|^2\le 1+(|z|+1/2)^2
                    \le \tfrac53(1+|z|^2).
\]
The second inequality follows by completing the square
in \( (2/3)|z|^2-|z|+5/12\).
Integrate the weight \(\prod_\tau(1+|u_\tau|^2)^{-2}\)
over each polydisk. Its volume is \((\pi/4)^f\)
and the integral of that weight over all
\(\mathbb C^f\) is \(\pi^f\).
The result follows first for any finite subset,
then for the full sum.

For a composition \(s_1+\cdots+s_b=s\), the nonzero
block-end minors \(D_1,\ldots,D_b\) determine the
block scalars up to at most \(\prod_h s_h\) choices:
successive quotients \(D_h/D_{h-1}\), up to the
fixed antidiagonal sign, prescribe their \(s_h\)-th
powers. In the original weight there is precisely
one factor \(W(D_h)\) for each block end.
We may sum over all integral \(D_h\), dropping every
further compatibility condition. The bound for this
composition is \(\prod_h s_h\,12^{fb}\).
Since \(s_h\le2^{s_h-1}\) for positive integers,
summing over the \(\binom{s-1}{b-1}\) compositions
of length \(b\) gives
\(12^f(12^f+2)^{s-1}\le14^{fs}\).
This proves (RM2) without a choice of integral
basis or a field-dependent packing constant.
\end{proof}

\begin{theorem}[Reduced actual generic point moment]
\label{thm:flag-reduced-moment}
Let \(J_2^B(w)\) be the original pinned squared orbital
amplitude, and use the original \(\mathrm{nPt},N_G,N_D,H\)
and the normalization
\[
 {\cal A}=AA\,\rho^{2r(r+1)f(k-1)}
\]
of (GP4). Put \(C=C_{1/2}\) from (NG4),
\(Q=|\mu(E)|^2 C^{2f}(TT')^{f/2}\), and
\(p=3s^2+s+2\). In the concrete equal-word case \(T'=T\).
Define the nonnegative point sum on occupied generic points
\[
 {\cal I}_p=\sum_w
 \mathrm{blockMinorWeight}(w)\,\mathrm{nPt}(w)\,
                      (N_G(w)N_D(w))^p J_2^B(w).
\]
Then the actual scalar corank-one generic portion satisfies
\[
 X.\mathrm{pairSum}(X.\mathrm{sc}\,X.\mathrm{genSet})
       \le 2 X.\mathrm{Gam1}\,{\cal A}\,
                         14^{fs}Q\,{\cal I}_p.           \tag{RM3}
\]
The inequality holds also with \(H(w)^p\) replacing
\((N_G(w)N_D(w))^p\) in the sum.
In particular the exponent \(6s^2\) in the prescribed
\(\mathrm{PointSumInput}\) suffices for this concrete
generic assembly. The general abstract
\(\mathrm{GenCutC}\) statement is not being asserted with
its original independent inputs and that smaller exponent.
\end{theorem}
\begin{proof}
Fix an occupied \(w\). Write \(V_{c,w,\nu}\)
for its actual grouped vectors \(\mathrm{grpGC}\), where
\(c\) is the scalar corank-one class and \(\nu\) is the
rational normal label. Use the finite witness of
Theorem~\ref{thm:flag-generic-pinned-input}
before its final enlargement of \(V_0\).
For any finite collection of these keys, (GP2) and
the exact row/datum weight comparison (GC4)--(GC5)
give
\[
 \sum_{c,\nu}\|V_{c,w,\nu}\|^2
 \le J_2^L(w) V_0
       \prod_{v\in P_*}
         \sum_{\delta_v}\mathrm{wt}_v(\delta_v)
                                  |A_v^0(\delta_v)|^2 .
\]
Apply (RM1) with the reference pairing classes of the
chosen generic term. Since \(H_v\ge2\) at every bad
place, \(2^{|P_*|}\le H(w)=\prod_{v\in P_*}H_v\).
Use the uninflated generic gluing mass (AN16),
\[
  V_0\prod_{v\in P_*}ngen_v
     \le 2{\cal A} H(w)^{s(s-1)}\mathrm{nPt}(w),
\]
and \(J_2^L(w)\le J_2^B(w)\). This proves
\[
 \sum_{c,\nu}\|V_{c,w,\nu}\|^2
 \le 2{\cal A} H(w)^{3s^2+s+1}
                       \mathrm{nPt}(w)J_2^B(w).        \tag{RM4}
\]
All weights and sums here are the original ones.
The right side is independent of the finite test
set, so the bound passes to the nonnegative sum.

At each fixed class and point (NG4) with \(\epsilon=1/2\)
bounds the number of occupied normal labels by
\(QG_0(w)\), where \(G_0=g^{-2sf}N_G\ge1\).
All the block-end minors at occupied points are integral
by (NG1). Put
\(W_*(w)=\mathrm{blockMinorWeight}(w)G_0(w)>0\).
For any finite collection of occupied groups in one
class, weighted Cauchy--Schwarz and (RM2) give
\[
 \left\|\sum_{w,\nu} V_{c,w,\nu}\right\|^2
 \le \left(\sum_{w,\nu} W_*(w)^{-1}\right)
       \sum_{w,\nu}W_*(w)\|V_{c,w,\nu}\|^2
 \le Q14^{fs}\sum_{w,\nu}W_*(w)\|V_{c,w,\nu}\|^2 .
\]
The original class-square expansion identifies
\(\mathrm{pairSum}(\mathrm{sc}\,\mathrm{genSet})\)
with \(\mathrm{Gam1}\) times the sum of these
class squares. This is precisely the
\(\mathrm{ClassSquaresSum}/\mathrm{RowFromClasses}\)
expansion; it uses scalar corank-one equivalence,
not an orthogonality assertion between unrelated
point groups. Absolute convergence of the original
expansion permits the grouping and the limits,
or the inequalities are obtained on finite groups
and then passed to the convergent class squares.
Summing the displayed Cauchy--Schwarz estimate over
classes and applying (RM4) proves the asserted bound
because the exact height formula is
\[
 H(w)=g^{-(2r+1)f}N_G(w)N_D(w)
       =g^fG_0(w)N_D(w).
\]
Both \(N_D\ge1\) and \(g\ge1\), so \(G_0\le H\)
and \(H\le N_GN_D\). In particular the original level
factor gives no additional positive power of \(g\)
in (RM3). Finally \(3s^2+s+2\le6s^2\) for \(s\ge1\)
and \(N_GN_D\ge1\).
\end{proof}

Thus the remaining real weighted estimate needed by
this route is the original pinned point sum with one
fixed weight exponent \(6s^2\), in place of the
\(43s^2\) obtained by the generic abstract assembly
with (GP3). The retained \(\mathrm{NormInputs}\) and
\(\mathrm{FrameInputs}\) normalize
\(\mathrm{Gam1}\,{\cal A}\,I_{\rm re}\) in precisely the
same way as in the source row estimate.
No claim is made here that the requisite
\({\cal I}_{6s^2}\) bound by
\(\exp(C_B(2s)^5 f\log b)I_{\rm re}\) has been proved.
In particular the equal-type real orbital estimate
on the singular block cells remains an analytic
problem even after the improvement (RM3).
